A registered paired protocol for a diffusion-personalisation comparison was
executed once. On the registered primary endpoint, per-pair identity fidelity,
the composed arm was lower on \RunIdentityLower{} of \RunPairs{} pairs and
higher on \RunIdentityHigher{} (sign-test \(p=\RunIdentitySignP{}\)); the
paired-difference interval \RunIdentityCiLo{} to \RunIdentityCiHi{} includes
zero, the pre-declared kill rule fired, and the test does not separate the
arms.

The protocol's size came from a recovered pairing. A recorded comparison had collapsed
\NPairs{} matched prompts into one dispersion per arm. Restored as pairs, the
composed arm was lower on identity for \IdentityDown{} of \NPairs{} prompts
(sign-test \(p=\IdentitySignP{}\)), and the smallest probability that test can
return at that size is \FloorAtFour{}, a design limit. The registered design
used \DesignTotal{} pairs across \DesignConcepts{} concepts, declared before
training (digest \PredeclarationShaShort{}) and realised with \RunFailures{}
failures.

Prompt fidelity, the registered secondary endpoint, was lower for the composed
arm on \RunPromptLower{} of \RunPairs{} pairs (\(p=\RunPromptSignP{}\)) and
stays secondary. Dispersion ratios (\RunDispersionDog{},
\RunDispersionClock{}) reversed the historical \SDRatioComposed{} and are
descriptive. The licensed claim is that a pre-declared paired sign test at this
size does not separate the arms on identity fidelity.
